% uid: ""
% school: "lycee"
% level: "3eme"
% section: ["Economie et Gestion"]
% chapter: "Suites réelles"
% title: "Dépréciation d'une voiture"
% difficulty: "Facile"
% points: null
% tags: ["suites", "suite-geometrique", "depreciation", "limite", "seuil"]
% provenance: "own"
% dressing: "suite géométrique décroissante"
% rung: 2
% batch: "2026-11-20-suites-3eco-A"

% ---STATEMENT---
Une voiture neuve vaut $15\,000$ dinars et perd $12\,\%$ de sa valeur chaque année. On note $U_n$ sa valeur au bout de $n$ années.

\begin{enumerate}
  \item Calculer $U_1$ puis montrer que $(U_n)$ est géométrique de raison $0{,}88$.
  \item Exprimer $U_n$ en fonction de $n$.
  \item Calculer $U_8$ et $U_9$ (arrondis au dixième). À partir de quelle année la voiture vaut-elle moins de $5\,000$ dinars ?
  \item Déterminer $\displaystyle\lim_{n\to+\infty}U_n$ et interpréter.
\end{enumerate}

% ---SOLUTION---
\begin{enumerate}
  \item $U_1=15\,000\times(1-0{,}12)=15\,000\times0{,}88=13\,200$ dinars.

  Perdre $12\,\%$ revient à multiplier par $1-0{,}12=0{,}88$. Donc pour tout $n$, $U_{n+1}=0{,}88\,U_n$ : $(U_n)$ est géométrique de raison $q=0{,}88$ et de premier terme $U_0=15\,000$.
  \item \textbf{Rappel.}\par
  \begin{prop}[Terme général d'une suite géométrique]
  $u_n=u_0\,q^n$.
  \end{prop}
  $$ \ib{$U_n=15\,000\times0{,}88^n$}.$$
  \item À la calculatrice :
  $$U_8=15\,000\times0{,}88^8\approx5\,394{,}5\qquad U_9=15\,000\times0{,}88^9\approx4\,747{,}2.$$
  $U_8>5\,000>U_9$ et la suite est décroissante : la voiture vaut moins de $5\,000$ dinars à partir de la $9^{\text{e}}$ année.
  \item $0<0{,}88<1$, donc $0{,}88^n\to0$ et $U_n\to15\,000\times0=$ \ib{$0$}.

  La valeur de la voiture se rapproche de $0$ sans jamais l'atteindre : la dépréciation est de plus en plus lente en dinars (on perd $12\,\%$ d'une valeur de plus en plus petite).
\end{enumerate}

\textbf{Idée :} pourcentage constant chaque année : on multiplie toujours par le même coefficient, c'est une suite géométrique.
